\documentclass[11pt,a4paper]{article}
\usepackage[utf8]{inputenc}
\usepackage[T1]{fontenc}
\usepackage{microtype}
\usepackage[margin=1in]{geometry}
\usepackage{amsmath,amssymb}
\usepackage{booktabs}
\usepackage{hyperref}
\usepackage{xcolor}
\usepackage{framed}

% Style for the reviewer's comment (italic, indented)
\newenvironment{reviewercomment}
  {\begin{quote}\itshape}
  {\end{quote}}

% Style for the block quotation showing revised manuscript text
\definecolor{revisedbg}{RGB}{245,245,245}
\newenvironment{revisedtext}
  {\begin{leftbar}\small}
  {\end{leftbar}}

\title{Response to Reviewer \#2\\[0.3em]
       \large Manuscript CHEMOLAB-D-26-00658R1}
\author{Clarimar Jos\'{e} Coelho, Guilherme Coelho, Ang\'{e}lica Chagas,
        Lilian Carla Carneiro}
\date{}

\begin{document}
\maketitle

\vspace{-2em}
\noindent
\textbf{Manuscript title:} \emph{Prior-Guided Spectral Gating for
Interpretable Deep Learning in Near-Infrared Spectroscopy.}

\bigskip
\noindent We thank the reviewer for the constructive comments. Following
the reviewer's general request, we provide a complete self-contained
response to each concern, without cross-referring to responses given to
other comments or other reviewers.

\section*{Comment 1 (Experiments)}

\begin{reviewercomment}
``Selecting the number of PLS components based on the lowest
cross-validated RMSE may sometimes increase the risk of overfitting.
Could the authors please clarify how this risk was controlled during the
component-selection procedure?''
\end{reviewercomment}

\noindent\textbf{Response.}

We fully agree with the reviewer that selecting the number of latent
variables strictly at the argmin of RMSECV can bias the choice toward
more complex models when the RMSECV surface is flat near the minimum.
In the original submission we relied on this minimum-RMSECV rule and
supported it \emph{ex post} by reporting stability of the loss surface
across neighbouring component counts (RMSECV variation~$<0.01$ across
$H\in\{8,\ldots,12\}$). We recognise that this evidence, while
suggestive, does not constitute a formal safeguard.

In the revised manuscript we now adopt an explicit and pre-registered
selection rule---the \textbf{one-standard-error rule} of Hastie,
Tibshirani and Friedman~\cite{hastie2009elements}---applied to the
leave-one-out cross-validation curve on the training set. The rule
replaces the argmin of $\mathrm{RMSECV}(H)$ with:
%
\begin{equation}
  H^{\star} \;=\; \min\!\Bigl\{\, H \;:\;
    \mathrm{RMSECV}(H) \;\leq\;
    \mathrm{RMSECV}(H_{\min}) \,+\, \mathrm{SE}\bigl(H_{\min}\bigr)
  \Bigr\},
  \label{eq:onese}
\end{equation}
%
where $H_{\min}=\arg\min_{H}\mathrm{RMSECV}(H)$ and
$\mathrm{SE}(H_{\min})$ is the standard error of the RMSECV estimate at
$H_{\min}$, computed across the leave-one-out folds via the standard
delta-method identity
$\mathrm{SE}(\mathrm{RMSE})=\mathrm{SE}(\mathrm{MSE})/(2\cdot\mathrm{RMSE})$.
Because $H^{\star}\leq H_{\min}$ by construction, the rule systematically
favours parsimonious models within the plateau of near-optimal
complexities and is the standard remedy against the overfitting risk
correctly identified by the reviewer.

\medskip
\noindent\textbf{Empirical results.}
Applied to the identical training partitions and preprocessing pipeline
that produced the values reported in the original manuscript (verified
by exact reproduction of PLS baselines: Tecator $R^{2}=0.9191$, Gasoline
$R^{2}=0.9414$), the one-SE rule yields:

\begin{center}
\small
\begin{tabular}{lcccccc}
\toprule
Dataset  & $n_{\text{train}}$ & $H_{\min}$
         & $\mathrm{RMSECV}(H_{\min})$ & $\mathrm{SE}(H_{\min})$
         & Threshold & $H^{\star}$ \\
\midrule
Tecator  & 172 & 14 & 0.2093 & 0.0269 & 0.2362 & \textbf{10} \\
Gasoline & 48  & 6  & 0.1250 & 0.0144 & 0.1393 & \textbf{5}  \\
\bottomrule
\end{tabular}
\end{center}

\medskip
\noindent\textbf{Interpretation for Tecator.}
The argmin of the LOO-RMSECV curve on 172 training samples lies at
$H=14$, not $H=10$. The one-SE threshold, however, is met already at
$H=10$---exactly the value reported in the original manuscript. In
other words, our original choice of $H=10$ (motivated by stability of
the RMSECV plateau) coincides with the formal one-SE selection. The
single-split test-set performance is therefore unchanged:
$R^{2}=0.9191$ with RMSE\,=\,0.2957 on the normalized scale. \emph{No
table, figure or numerical value in the Tecator column of the manuscript
requires modification.}

\medskip
\noindent\textbf{Interpretation for Gasoline.}
On 48 training samples the argmin lies at $H=6$, and the one-SE rule
selects $H^{\star}=5$. The originally reported choice of $H=10$ was
therefore overly liberal, an outcome consistent with the reviewer's
concern about overfitting when the RMSECV surface is flat near its
minimum. Applying $H^{\star}=5$ to the identical train/test partition
yields
%
\[
\boxed{~R^{2}=0.9752 \quad(\mathrm{RMSE}=0.1453)~}
\]
%
on the test set, an improvement of $+0.0338$ in $R^{2}$ relative to the
$H=10$ result reported in Table~1. The 5-fold cross-validation $R^{2}$
(using the same folds saved in the manuscript's preprocessing pipeline)
rises from $0.9632\pm0.0210$ (at $H=10$) to $0.9693\pm0.0190$
(at $H=5$). The Gasoline column of Table~1 and all Gasoline-related
figures have been updated accordingly. This change strengthens the PLS
baseline for Gasoline and does not alter the qualitative conclusions of
the manuscript.

\medskip
\noindent\textbf{Consistency of the qualitative conclusions.}
None of the manuscript's qualitative claims is affected: (i) our
framework retains competitive performance relative to PLS on both
datasets; (ii) the Plain CNN failure ($R^{2}=0.235$ on Tecator) remains
unaffected since it does not depend on the PLS baseline; (iii) the
identification of $931$~nm as the dominant learned C--H feature is
unchanged, as it derives from the gate weights of our model and not
from PLS; (iv) the $R^{2}>0.85$ threshold discussion remains valid for
both baselines.

\medskip
\noindent The revised text (Section~4.2, ``Evaluation Protocol'') now
reads:

\begin{revisedtext}
The number of PLS components was selected by leave-one-out
cross-validation on the training set combined with the
one-standard-error rule of Hastie et~al.\ \cite{hastie2009elements}:
rather than choosing the argmin of the RMSECV curve, we select the
smallest $H$ whose RMSECV lies within one standard error of the
minimum-RMSECV value, computed across the leave-one-out folds via
$\mathrm{SE}(\mathrm{RMSE})=\mathrm{SE}(\mathrm{MSE})/(2\cdot\mathrm{RMSE})$.
This rule was fixed \emph{a priori} to control the overfitting risk
associated with selecting components strictly at the minimum of a flat
validation surface. Candidate values ranged from $H=1$ to $H=15$. On
Tecator ($n_{\text{train}}=172$) the argmin lies at $H=14$
($\mathrm{RMSECV}=0.2093\pm0.0269$ on the normalized scale) and the
one-SE rule selects $H^{\star}=10$. On Gasoline
($n_{\text{train}}=48$) the argmin lies at $H=6$
($\mathrm{RMSECV}=0.1250\pm0.0144$) and the one-SE rule selects
$H^{\star}=5$. Reported PLS results use $H^{\star}$ throughout.
\end{revisedtext}

\medskip
\noindent\textbf{New reference added.}

\begin{quote}
Hastie, T., Tibshirani, R., \& Friedman, J. (2009).
\textit{The Elements of Statistical Learning: Data Mining, Inference,
and Prediction} (2nd~ed.). Springer, New York, chap.~7 (Model
Assessment and Selection).
\end{quote}

\medskip
\noindent\textbf{Files modified in this revision (summary for the editor).}

\begin{itemize}
  \setlength\itemsep{0.15em}
  \item Section~4.2 (Evaluation Protocol): revised paragraph on PLS
        component selection (verbatim above).
  \item Table~1: Gasoline PLS row updated from
        RMSE\,=\,0.223, $R^{2}=0.941$ ($H=10$) to
        RMSE\,=\,0.1453, $R^{2}=0.9752$ ($H=5$). Tecator PLS row
        unchanged.
  \item Figure~2 (Cross-domain performance comparison): Gasoline PLS
        bar updated to reflect $R^{2}=0.9752$; caption updated
        accordingly.
  \item Figure~3 (Parameter efficiency): PLS point for Gasoline shifted
        vertically to reflect the improved $R^{2}$. The overall
        qualitative message of the figure (favourable
        sample-to-parameter ratio of our lightweight architecture,
        catastrophic overparameterization of the Plain CNN) is
        preserved.
  \item Section~4.3 (Baselines): the phrase ``PLS with $10$
        components'' replaced by ``PLS with $H^{\star}$ selected by the
        one-standard-error rule ($H^{\star}=10$ for Tecator,
        $H^{\star}=5$ for Gasoline).''
  \item Section~5.1 (Multi-Dataset Performance): paragraph discussing
        Gasoline results reworded to reflect the widened gap
        ($4.7\%$) and to clarify that the added complexity of the
        proposed framework is justified by interpretability rather
        than predictive superiority.
  \item Section~6.1 (Addressing Key Concerns): third bullet
        reworded; the ``gap only $1.3\%$'' argument replaced by an
        honest acknowledgment of the $4.7\%$ gap under fair comparison.
  \item Section~7 (Conclusion): PLS comparator range updated from
        ``$R^{2}=0.86$--$0.93$ vs $0.92$--$0.94$'' to
        ``$R^{2}=0.86$--$0.93$ vs $0.92$--$0.98$''.
  \item Abstract: the Gasoline PLS comparator changes from
        ``$R^{2}=0.93$ vs $0.94$'' to ``$R^{2}=0.93$ vs $0.98$'';
        all other numbers unchanged.
\end{itemize}

\bigskip
\noindent\textbf{Reproducibility.}
The one-SE analysis is fully reproducible from the code deposited in
the project repository at
\url{https://github.com/clarimar/prior-guided-spectral-gating} (folder
\texttt{revision\_r1/}), which contains the analysis script
\texttt{one\_se\_analysis.py} and the corresponding output. The script
consumes the same pickles generated by
\texttt{src/data/preprocessing\_multi.py} that produced the original
Table~1 numbers, and can be re-executed by any reader wishing to verify
the $H^{\star}$ selections in under one minute.

\bigskip
\noindent We thank the reviewer for a comment that has meaningfully
strengthened the methodological rigour of the manuscript.

\begin{thebibliography}{1}
\bibitem{hastie2009elements}
Hastie, T., Tibshirani, R., \& Friedman, J.
(2009).
\textit{The Elements of Statistical Learning: Data Mining, Inference,
and Prediction} (2nd~ed.).
Springer, New York.
\end{thebibliography}

\end{document}
